% ==============================================================================
% RESUMO E ABSTRACT
% ==============================================================================
\newpage
\begin{center}
    \textbf{RESUMO}
\end{center}
\vspace{0.3cm}

\noindent A rápida expansão de modelos de inteligência artificial generativa e redes neurais profundas intensificou a crise de sustentabilidade da infraestrutura computacional terrestre, caracterizada pela sobrecarga de redes elétricas metropolitanas e pelo consumo de bilhões de litros de água em resfriamento evaporativo. Como alternativa emergente, propostas de vanguarda aeroespacial — notadamente o consórcio europeu ASCEND, o projeto Starmind da SpaceX e os estudos de infraestrutura escalável de IA no espaço desenvolvidos por pesquisadores do Google (AGÜERA Y ARCAS et al., 2025; BLOCH et al., 2026) — vislumbram a instalação de Centros de Dados Orbitais (\textit{Orbital Data Centers} -- ODCs) em órbitas com alta ou contínua insolação solar (como órbitas heliossíncronas \textit{dawn-dusk} e órbitas geoestacionárias). Sob o espectro extraterrestre AM0 ($1361\text{ W/m}^2$), a geração fotovoltaica é contínua e até oito vezes mais produtiva que em solo, desonerando redes elétricas e preservando recursos hídricos. Contudo, essa operação revela um gargalo físico limitante: o gerenciamento térmico no vácuo espacial. Na ausência de convecção ou condução com o ambiente circundante, toda a expressiva potência térmica dissipada por processadores de IA (GPUs, TPUs e aceleradores neurais) deve ser rejeitada exclusivamente por radiação térmica regida pela Lei de Stefan-Boltzmann ($q = \varepsilon \sigma A (T_{rad}^4 - T_{space}^4)$). Sem gerenciamento ativo, a temperatura de junção dos chips atinge rapidamente os limites de tolerância ($T_j \le 85^\circ\text{C}$), ao passo que o sobredimensionamento de radiadores térmicos acarreta penalidades proibitivas de massa e volume de lançamento. Esta dissertação propõe o desenvolvimento, a parametrização e a validação de um framework computacional integrado de modelagem matemática e simulação para o gerenciamento térmico-energético de data centers orbitais. A pesquisa integra: (i) modelagem fotovoltaica espacial de alta potência sob espectro AM0 acoplada a algoritmos de rastreamento do ponto de máxima potência (MPPT); (ii) caracterização eletrotérmica de estágios de distribuição de potência DC-DC de alta eficiência; (iii) modelagem de baterias para sustentação transitória supervisionada por Filtro de Kalman Unscented adaptativo (UKF); (iv) equacionamento do balanço térmico radiativo a vácuo para nós de processamento e radiadores espaciais; e (v) algoritmo de Sistema de Gerenciamento Térmico-Energético (TEMS) com escalonamento térmico-consciente de tarefas de IA e modulação dinâmica de tensão e frequência (DVFS). Adotando metodologia quantitativa estritamente computacional em MATLAB/Simulink e Python, a pesquisa fornece à engenharia de controle e à indústria aeroespacial uma plataforma preditiva aberta para projeto, controle ótimo e operação segura de servidores de IA no espaço.

\vspace{0.5cm}
\noindent\textbf{Palavras-chave:} Data Centers Orbitais; Inteligência Artificial Espacial; Gerenciamento Térmico no Vácuo; Sistema de Energia Elétrica (EPS); Escalonamento Térmico-Consciente; Lei de Stefan-Boltzmann; Modelagem Computacional.

\newpage
\begin{center}
    \textbf{ABSTRACT}
\end{center}
\vspace{0.3cm}

\noindent The rapid expansion of generative artificial intelligence models and deep neural networks has intensified the sustainability crisis of terrestrial computing infrastructure, marked by the overburdening of metropolitan power grids and the consumption of billions of liters of water for evaporative cooling. As an emerging alternative, leading aerospace initiatives — notably the European ASCEND consortium, SpaceX's Starmind project, and scalable space AI infrastructure research from Google (AGÜERA Y ARCAS et al., 2025; BLOCH et al., 2026) — envision deploying Orbital Data Centers (ODCs) into high or continuous insolation orbits (such as dawn-dusk Sun-Synchronous Orbits and Geostationary Earth Orbits). Under the extraterrestrial AM0 solar spectrum ($1361\text{ W/m}^2$), photovoltaic power generation is continuous and up to eight times more productive than on Earth's surface, offloading terrestrial electrical grids and preserving water resources. However, this operation exposes a severe physical bottleneck: thermal management in space vacuum. In the absence of convective or conductive heat transfer with the surrounding environment, the massive thermal power dissipated by AI hardware accelerators (GPUs, TPUs, and neural engines) must be rejected strictly via surface thermal radiation governed by the Stefan-Boltzmann law ($q = \varepsilon \sigma A (T_{rad}^4 - T_{space}^4)$). Without active management, semiconductor junction temperatures rapidly exceed operational safety limits ($T_j \le 85^\circ\text{C}$), whereas oversizing deployable space radiators imposes prohibitive launch mass and volume penalties. This Master's thesis proposes the development, parameterization, and computational validation of an integrated mathematical modeling and dynamic simulation framework for thermal-energy management of orbital data centers. The research integrates: (i) high-power space photovoltaic array modeling under the AM0 spectrum coupled with adaptive Maximum Power Point Tracking (MPPT) algorithms; (ii) electrothermal characterization of high-efficiency DC-DC power distribution stages; (iii) lithium-ion battery energy storage modeling for transient stabilization supervised by an adaptive Unscented Kalman Filter (UKF) for State of Charge (SoC) estimation; (iv) vacuum radiative heat balance formulation coupling processor nodes and space radiators; and (v) a Thermal-Energy Management System (TEMS) algorithm incorporating thermal-aware AI workload scheduling and Dynamic Voltage and Frequency Scaling (DVFS). Adopting a strictly theoretical-computational methodology implemented in MATLAB/Simulink and Python, this research provides the control engineering community and the aerospace industry with an open predictive simulation platform for design, optimal control, and reliable operation of AI servers in space.

\vspace{0.5cm}
\noindent\textbf{Keywords:} Orbital Data Centers; Space Artificial Intelligence; Space Vacuum Thermal Management; Electrical Power System (EPS); Thermal-Aware Scheduling; Stefan-Boltzmann Law; Computational Modeling.